\subsection{\texorpdfstring{SYNODALITY AND PRIMACY DURING THE FIRST
MILLENNIUM:\\
TOWARDS A COMMON UNDERSTANDING\\
IN SERVICE TO THE UNITY OF THE
CHURCH}{SYNODALITY AND PRIMACY DURING THE FIRST MILLENNIUM: TOWARDS A COMMON UNDERSTANDING IN SERVICE TO THE UNITY OF THE CHURCH}}\label{synodality-and-primacy-during-the-first-millennium-towards-a-common-understanding-in-service-to-the-unity-of-the-church}

\subsubsection{Chieti, 21 September
2016}\label{chieti-21-september-2016}

~

\emph{We declare to you what we have seen and heard so that you also may
have communion {[}koinonia{]} with us;\\
and truly our communion {[}koinonia{]} is with the Father and with his
Son Jesus Christ.\\
We are writing these things so that our joy may be complete.'\\
(1 Jn 1:3-4)}

~

1. Ecclesial communion arises directly from the Incarnation of the
eternal Word of God, according to the goodwill (\emph{eudokia}) of the
Father, through the Holy Spirit. Christ, having come on earth, founded
the Church as his body (cf. 1Cor 12:12-27). The unity that exists among
the Persons of the Trinity is reflected in the communion
(\emph{koinonia}) of the members of the Church with one another. Thus,
as St Maximus the Confessor affirmed, the Church is an~\emph{`eikon'~}of
the Holy Trinity.(1) At the Last Supper, Jesus Christ prayed to his
Father: `Protect them in your name that you have given me, so that they
may be one, as we are one' (Jn 17:11). This Trinitarian unity is
manifested in the Holy Eucharist, wherein the Church prays to God the
Father through Jesus Christ in the Holy Spirit.

2. From earliest times, the one Church existed as many local churches.
The communion (\emph{koinonia}) of the Holy Spirit (cf. 2Cor 13:13) was
experienced both within each local church and in the relations between
them as a unity in diversity. Under the guidance of the Spirit (cf. Jn
16:13), the Church developed patterns of order and various practices in
accordance with its nature as `a people brought into unity from the
unity of the Father, the Son and the Holy Spirit'.(2)

3. Synodality is a fundamental quality of the Church as a whole. As St
John Chrysostom said: `"Church" means both gathering
{[}\emph{systema}{]} and synod {[}\emph{synodos}{]}'.(3)~The term comes
from the word `council' (\emph{synodos~}in Greek,~\emph{concilium~}in
Latin), which primarily denotes a gathering of bishops, under the
guidance of the Holy Spirit, for common deliberation and action in
caring for the Church. Broadly, it refers to the active participation of
all the faithful in the life and mission of the Church.

4. The term primacy refers to being the first
(\emph{primus},~\emph{protos}). In the Church, primacy belongs to her
Head -- Jesus Christ, `who is the beginning, the firstborn from the
dead; that in all things he might have the pre-eminence
{[}\emph{protevon}{]}' (Col. 1:18). Christian Tradition makes it clear
that, within the synodal life of the Church at various levels, a bishop
has been acknowledged as the `first'. Jesus Christ associates this being
`first' with service (\emph{diakonia}): `Whoever wants to be first must
be last of all and servant of all' (Mk 9:35).

5. In the second millennium, communion was broken between East and West.
Many efforts have been made to restore communion between Catholics and
Orthodox, but they have not succeeded. The Joint International
Commission for Theological Dialogue between the Roman Catholic Church
and the Orthodox Church, in its ongoing work to overcome theological
divergences, has been considering the relationship between synodality
and primacy in the life of the Church. Different understandings of these
realities played a significant role in the division between Orthodox and
Catholics. It is, therefore, essential to seek to establish a common
understanding of these interrelated, complementary and inseparable
realities.

6. In order to achieve this common understanding of primacy and
synodality, it is necessary to reflect upon history. God reveals himself
in history. It is particularly important to undertake together a
theological reading of the history of the Church's liturgy,
spirituality, institutions and canons, which always have a theological
dimension.

7. The history of the Church in the first millennium is decisive.
Despite certain temporary ruptures, Christians from East and West lived
in communion during that time, and, within that context, the essential
structures of the Church were constituted. The relationship between
synodality and primacy took various forms, which can give vital guidance
to Orthodox and Catholics in their efforts to restore full communion
today.

\subsubsection{The Local Church}\label{the-local-church}

8. The one, holy, catholic and apostolic Church of which Christ is the
head is present in the eucharistic synaxis of a local church under its
bishop. He is the one who presides (the~\emph{`proestos'}). In the
liturgical synaxis, the bishop makes visible the presence of Jesus
Christ. In the local church (i.e. a diocese), the many faithful and
clergy under the one bishop are united with one another in Christ, and
are in communion with him in every aspect of the life of the Church,
most especially in the celebration of the Eucharist. As St Ignatius of
Antioch taught: `where the bishop is, there let all the people be, just
as, where Jesus Christ is, we have the catholic church
{[}\emph{katholike ekklesia}{]}'.(4)~Each local church celebrates in
communion with all other local churches which confess the true faith and
celebrate the same Eucharist. When a presbyter presides at the
Eucharist, the local bishop is always commemorated as a sign of the
unity of the local church. In the Eucharist, the~\emph{proestos~}and the
community are interdependent: the community cannot celebrate the
Eucharist without a~\emph{proestos}, and the~\emph{proestos}, in turn,
must celebrate with a community.

9. This interrelatedness between the~\emph{proestos~}or bishop and the
community is a constitutive element of the life of the local church.
Together with the clergy, who are associated with his ministry, the
local bishop acts in the midst of the faithful, who are Christ's flock,
as guarantor and servant of unity. As successor of the Apostles, he
exercises his mission as one of service and love, shepherding his
community, and leading it, as its head, to ever-deeper unity with Christ
in the truth, maintaining the apostolic faith through the preaching of
the Gospel and the celebration of the sacraments.

10. Since the bishop is the head of his local church, he represents his
church to other local churches and in the communion of all the churches.
Likewise, he makes that communion present to his own church. This is a
fundamental principle of synodality.

\subsubsection{The Regional Communion of
Churches}\label{the-regional-communion-of-churches}

11. There is abundant evidence that bishops in the early Church were
conscious of having a shared responsibility for the Church as a whole.
As St Cyprian said: `There is but one episcopate but it is spread
amongst the harmonious host of all the numerous bishops'.(5)~This bond
of unity was expressed in the requirement that at least three bishops
should take part in the ordination (\emph{cheirotonia}) of a new
one;(6)it was also evident in the multiple gatherings of bishops in
councils or synods to discuss in common issues of doctrine
(\emph{dogma},~\emph{didaskalia}) and practice, and in their frequent
exchanges of letters and mutual visits.

12. Already during the first four centuries, various groupings of
dioceses within particular regions emerged. The~\emph{protos}, the first
among the bishops of the region, was the bishop of the first see, the
metropolis, and his office as metropolitan was always attached to his
see. The ecumenical councils attributed certain prerogatives
(\emph{presbeia},~\emph{pronomia},~\emph{dikaia}) to the metropolitan,
always within the framework of synodality. Thus, the First Ecumenical
Council (Nicaea, 325), while requiring of all the bishops of a province
their personal participation in or written agreement to an episcopal
election and consecration - a synodical act~\emph{par excellence~}-
attributed to the metropolitan the validation (\emph{kyros}) of the
election of a new bishop.(7) The Fourth Ecumenical Council (Chalcedon,
451) again evoked the rights (\emph{dikaia}) of the metropolitan --
insisting that this office is ecclesial, not political~(8)~- as did the
Seventh Ecumenical Council (Nicaea II, 787), also.(9)

13. Apostolic Canon 34 offers a canonical description of the correlation
between the~\emph{protos~}and the other bishops of each region: `The
bishops of the people of a province or region {[}\emph{ethnos}{]} must
recognize the one who is first {[}\emph{protos}{]} amongst them, and
consider him to be their head {[}\emph{kephale}{]}, and not do anything
important without his consent {[}\emph{gnome}{]}; each bishop may only
do what concerns his own diocese {[}\emph{paroikia}{]} and its dependent
territories. But the first {[}\emph{protos}{]} cannot do anything
without the consent of all. For in this way concord
{[}\emph{homonoia}{]} will prevail, and God will be praised through the
Lord in the Holy Spirit'.(10)

14. The institution of the metropolitanate is one form of regional
communion between local churches. Subsequently other forms developed,
namely the patriarchates comprising several metropolitanates. Both a
metropolitan and a patriarch were diocesan bishops with full episcopal
power within their own dioceses. In matters related to their respective
metropolitanates or patriarchates, however, they had to act in accord
with their fellow bishops. This way of acting is at the root of
synodical institutions in the strict sense of the term, such as a
regional synod of bishops. These synods were convened and presided over
by the metropolitan or the patriarch. He and all the bishops acted in
mutual complementarity and were accountable to the synod.

\subsubsection{The Church at the Universal
Level}\label{the-church-at-the-universal-level}

15. Between the fourth and the seventh centuries, the order
(\emph{taxis}) of the five patriarchal sees came to be recognised, based
on and sanctioned by the ecumenical councils, with the see of Rome
occupying the first place, exercising a primacy of honour
(\emph{presbeia tes times}), followed by the sees of Constantinople,
Alexandria, Antioch and Jerusalem, in that specific order, according to
the canonical tradition.(11)

16. In the West, the primacy of the see of Rome was understood,
particularly from the fourth century onwards, with reference to Peter's
role among the Apostles. The primacy of the bishop of Rome among the
bishops was gradually interpreted as a prerogative that was his because
he was successor of Peter, the first of the apostles.(12)~This
understanding was not adopted in the East, which had a different
interpretation of the Scriptures and the Fathers on this point. Our
dialogue may return to this matter in the future.

17. When a new patriarch was elected to one of the five sees in
the~\emph{taxis}, the normal practice was that he would send a letter to
all the other patriarchs, announcing his election and including a
profession of faith. Such `letters of communion' profoundly expressed
the canonical bond of communion among the patriarchs. By including the
new patriarch's name, in the proper order, in the diptychs of their
churches, read in the Liturgy, the other patriarchs acknowledged his
election. The~\emph{taxis~}of the patriarchal sees had its highest
expression in the celebration of the holy Eucharist. Whenever two or
more patriarchs gathered to celebrate the Eucharist, they would stand
according to the~\emph{taxis}. This practice manifested the eucharistic
character of their communion.

18. From the First Ecumenical Council (Nicaea, 325) onwards, major
questions regarding faith and canonical order in the Church were
discussed and resolved by the ecumenical councils. Though the bishop of
Rome was not personally present at any of those councils, in each case
either he was represented by his legates or he agreed with the council's
conclusions~\emph{post factum}. The Church's understanding of the
criteria for the reception of a council as ecumenical developed over the
course of the first millennium. For example, prompted by historical
circumstances, the Seventh Ecumenical Council (Nicaea II, 787) gave a
detailed description of the criteria as then understood: the agreement
(\emph{symphonia}) of the heads of the churches, the cooperation
(\emph{synergeia}) of the bishop of Rome, and the agreement of the other
patriarchs (\emph{symphronountes}). An ecumenical council must have its
own proper number in the sequence of ecumenical councils, and its
teaching must accord with that of previous councils.(13) Reception by
the Church as a whole has always been the ultimate criterion for the
ecumenicity of a council.

19. Over the centuries, a number of appeals were made to the bishop of
Rome, also from the East, in disciplinary matters, such as the
deposition of a bishop. An attempt was made at the Synod of Sardica
(343) to establish rules for such a procedure.(14)~Sardica was received
at the Council~\emph{in Trullo~}(692).(15)~The canons of Sardica
determined that a bishop who had been condemned could appeal to the
bishop of Rome, and that the latter, if he deemed it appropriate, might
order a retrial, to be conducted by the bishops in the province
neighbouring the bishop's own. Appeals regarding disciplinary matters
were also made to the see of Constantinople,(16)~and to other sees. Such
appeals to major sees were always treated in a synodical way. Appeals to
the bishop of Rome from the East expressed the communion of the Church,
but the bishop of Rome did not exercise canonical authority over the
churches of the East.

\subsubsection{Conclusion}\label{conclusion}

20. Throughout the first millennium, the Church in the East and the West
was united in preserving the apostolic faith, maintaining the apostolic
succession of bishops, developing structures of synodality inseparably
linked with primacy, and in an understanding of authority as a service
(\emph{diakonia}) of love. Though the unity of East and West was
troubled at times, the bishops of East and West were conscious of
belonging to the one Church.

21. This common heritage of theological principles, canonical provisions
and liturgical practices from the first millennium constitutes a
necessary reference point and a powerful source of inspiration for both
Catholics and Orthodox as they seek to heal the wound of their division
at the beginning of the third millennium. On the basis of this common
heritage, both must consider how primacy, synodality, and the
interrelatedness between them can be conceived and exercised today and
in the future.

~

\paragraph{ENDNOTES}\label{endnotes}

1) St Maximus the Confessor,~\emph{Mystagogia~}(PG 91, 663D).

2) St Cyprian,~\emph{De Orat. Dom}., 23 (PL 4, 536).

3) St John Chrysostom,~\emph{Explicatio in Ps 149~}(PG 55, 493).

4) St Ignatius,~\emph{Letter to the Smyrnaeans}, 8.

5) St Cyprian,~\emph{Ep.}55, 24, 2; cf. also,~\emph{`episcopatus unus
est cuius a singulis in solidum pars tenetur'~}(De unitate, 5).

6) First Ecumenical Council (Nicaea, 325), canon 4:~`It is preferable
that a bishop be established by all the bishops of a province; but if
this appears difficult because of a pressing necessity or because of the
distance to be travelled, at least three bishops should come together;
and, having the written consent of the absent bishops, they may then
proceed with the consecration. The validation {[}\emph{kyros}{]} of what
takes place falls on the metropolitan bishop of each province.' Cf.
also~\emph{Apostolic Canon}, 1: `A bishop must be ordained by two or
three bishops'.

7) First Ecumenical Council (Nicaea, 325), canon 4; also canon 6: `If
anyone becomes a bishop without the consent of the metropolitan, the
great council decrees that such a person is not even a bishop.'

8) Fourth Ecumenical Council (Chalcedon, 451), canon 12: `As for cities
that have already been honoured by the title of metropolis by imperial
letters, let these cities and the bishops who govern them enjoy only the
honour of the title; that is, let the proper rights of the true
{[}\emph{kata aletheian}{]} metropolis be safeguarded.'

9) Seventh Ecumenical Council (Nicaea II, 787), canon 11 grants the
metropolitans the right to appoint the treasurers of their suffragan
dioceses if the bishops do not provide for it.

10) Cf. Council of Antioch (327), canon 9: `It is proper for the bishops
in every province {[}\emph{eparchia}{]} to submit to the bishop who
presides in the metropolis'.

11) Cf. First Ecumenical Council (Nicaea, 325), canon 6: `The ancient
customs of Egypt, Libya and Pentapolis shall be maintained, according to
which the bishop of Alexandria has authority over all these places,
since a similar custom exists with reference to the bishop of Rome.
Similarly in Antioch and the other provinces, the prerogatives
{[}\emph{presbeia}{]} of the churches are to be preserved'; Second
Ecumenical Council (Constantinople, 381), canon 3: Let the bishop of
Constantinople \ldots{} have the primacy of honour {[}\emph{presbeia tes
times}{]} after the bishop of Rome, because it is New Rome'; Fourth
Ecumenical Council (Chalcedon, 451), canon 28: `The Fathers rightly
accorded prerogatives {[}\emph{presbeia}{]} to the see of older Rome
since that is an imperial city; and moved by the same purpose the one
hundred and fifty most devout bishops apportioned equal prerogatives to
the most holy see of New Rome, reasonably judging that the city which is
honoured by the imperial power and senate and enjoying privileges
equalling older imperial Rome, should also be elevated to her level in
ecclesiastical affairs and take second place after her' (this canon was
never received in the West); Council~\emph{in Trullo~}(692), canon 36:
`Renewing the enactments of the one hundred and fifty Fathers assembled
at the God-protected and imperial city, and those of the six hundred and
thirty who met at Chalcedon, we decree that the see of Constantinople
shall have equal privileges {[}\emph{presbeia}{]} with the see of Old
Rome, and shall be highly regarded in ecclesiastical matters as that see
is and shall be second after it. After Constantinople shall be ranked
the see of Alexandria, then that of Antioch, and afterwards the see of
Jerusalem'.

12) Cf. Jerome,~\emph{In Isaiam~}14, 53; Leo,~\emph{Sermo~}96, 2-3.

13) Cf. Seventh Ecumenical Council (Nicaea II, 787): J. D.
MANSI,~\emph{Sacrorum conciliorum nova et amplissima collectio}, XIII,
208D-209C.

14) Cf. Synod of Sardica (343), canons 3 and 5.

15) Cf. Council~\emph{in Trullo}, canon 2. Similarly, the Photian
Council of 861 accepted the canons of Sardica as recognising the bishop
of Rome as having a right of cassation in cases already judged in
Constantinople.

16) Cf. Fourth Ecumenical Council (Chalcedon, 451), canons 9 and 17.
